% TOP_PROBLEM: {"id": 5300009, "title": "Polynomial realization of intertwining", "category_id": 11, "category": {"id": 11, "name": "dynamical_systems", "display_name": "Dynamical Systems", "description": "Problems about long-term behavior of deterministic systems, Hamiltonian dynamics, and periodic orbits.", "slug": "dynamical-systems", "order_index": 11, "created_at": "2026-07-31T15:26:25.671Z"}, "status": "open"}
% FIRST_POSTED: null
% ENTRY_KIND: statement_only
% Imported from: batch04.tex; UnsolvedMath ID 5300009
\documentclass[11pt]{article}
\usepackage{fontspec}
\setmainfont{Latin Modern Roman}
\setsansfont{Latin Modern Sans}
\usepackage{amsmath,amssymb,mathtools}
\usepackage{unicode-math}
\setmathfont{Latin Modern Math}
\usepackage{xurl}
\usepackage[hidelinks]{hyperref}
\newcommand{\sref}[2]{\hyperlink{source:#1:#2}{[S#2]}}
\newcommand{\cataloguescope}[1]{\par\noindent\textbf{Question.} #1\par}
\begin{document}
% UnsolvedMath ID 5300009; source catalogue code AMR-052-0009
\setcounter{section}{5300008}
\section{Polynomial realization of intertwining}\label{5300009}
{\small\sffamily UnsolvedMath ID 5300009 · Dynamical Systems}
\subsection{Definitions and mathematical statement}

\paragraph{Shared notation.}$\mathbb N=\{1,2,\ldots\}$, and $\mathbb Z,\mathbb Q,\mathbb R,\mathbb C$ denote the integers, rationals, reals and complex numbers. Any other conventions are specified locally.


For a complex polynomial $P$ of degree at least two, its filled Julia set $K(P)$ is the set of points with bounded forward orbit, and $J(P)=\partial K(P)$ is its Julia set. An external ray is a radial curve in the coordinate near infinity that conjugates $P$ to its leading power map, continued into the basin of infinity. Its angle is measured in $\mathbb R/\mathbb Z$. A critical point is a zero of $P'$. Topological conjugacy means a homeomorphism $h$ satisfying $h\circ f=g\circ h$. A quasiconformal homeomorphism has bounded infinitesimal eccentricity; quasiconformal surgery cuts and glues dynamical planes and finds an invariant complex structure that can be straightened to a holomorphic map.
Let $P_1$ be monic with connected Julia set and a repelling fixed point $x_0$. Choose an external ray landing at $x_0$ whose forward cycle consists of $q$ rays with combinatorial rotation number $p/q$: the dynamics permutes the cyclically ordered rays by a shift of $p$ positions. These rays cut the first dynamical plane into $q$ wedges. Let $P_2$ be monic with a ray of the same combinatorial rotation number landing at a repelling periodic point whose period divides $q$. Slit the second plane along the corresponding rays and insert the wedges from the first plane in the corresponding order. On the glued surface use $P_1$ and $P_2$ away from the cuts, adjusting the dynamics near inverse images of the cut rays to obtain continuity. This is the intertwining construction. The second landing orbit is required to be repelling; the source explicitly distinguishes the parabolic case, where this construction does not seem to work.
\paragraph{Statement recorded in the supplied source.}
For which admissible intertwining constructions is the resulting map conjugate to a polynomial? \sref{5300009}{2}
\subsection{Short English statement}
When does intertwining the two polynomial planes produce polynomial dynamics?
\subsection{Sources}
\begin{description}
\item[\hypertarget{source:5300009:1}{S1}] ULAM.AI and the UnsolvedMath Contributors, UnsolvedMath: A Curated Collection of Open Mathematics Problems, record 5300009 (AMR-052-0009). CC BY 4.0. \url{https://huggingface.co/datasets/ulamai/UnsolvedMath}.
\item[\hypertarget{source:5300009:2}{S2}] Ben Bielefeld, Questions in Quasiconformal Surgery, in Ben Bielefeld and Mikhail Lyubich (editors), Problems in holomorphic dynamics (1992). Intertwining definition and Question 9, pp. 4–5. \url{https://arxiv.org/pdf/math/9205209}.
\end{description}
\label{end:5300009}
\end{document}
